% GENERATED FILE. Edit canonical lessons, then run npm run generate:books.
% canonical-source-hash: fnv1a64:4c993c018ca4d91f
% canonical-lessons: GU-C528-yadi, GU-C528-bajet, GU-C528-vyavastha, GU-C528-dayri, GU-C528-yojvu

\chapter{If We Plan Ahead}
\label{ch:gu-if-we-plan-ahead}
\hlchaptermodality{528}

\begin{chapteropening}
\textbf{By the end of this chapter:} \emph{I can say what will happen if we plan ahead.}

Complete the last lesson of chapter 528: I can say what will happen if we plan ahead.
\end{chapteropening}

\section[\texorpdfstring{\gu{યાદી} (yādī)}{યાદી (yādī)}]{\gu{યાદી} (yādī) — a list}

\label{lesson:GU-C528-yadi}



\begin{quote}
Before the new one: say the Gujarati for to throw down, then the Gujarati for to tire out (an opponent).
\end{quote}

\subsection*{You'll want to know: \gu{યાદી}}
\textbf{\gu{યાદી}} — \emph{yādī} — \textquotedblleft{}a list\textquotedblright{}.

\begin{quote}
\gu{જો} \gu{તમે} \gu{યાદી} \gu{બનાવશો} \gu{તો} \gu{તમે} \gu{કશું} \gu{નહીં} \gu{ભૂલો}.
\end{quote}

\begin{quote}
\emph{jo tame yādī banāvśo to tame kaśũ nahī̃ bhūlo.}
\end{quote}

\begin{quote}
If you make a list, you won't forget anything.
\end{quote}

A feminine noun.

\begin{cousinweb}
\textbf{\gu{યાદી}} (\emph{yādī}) is usually traced to Persian \emph{yād}, \textquotedblleft{}memory\textquotedblright{}, with the ending \emph{-ī}: something written down to remember by.
\end{cousinweb}

\begin{grammarlens}[title={What you've built}]
One more word for saying what happens if.
\end{grammarlens}

\subsection*{Your turn}
\begin{itemize}
\raggedright
  \item \emph{Say it:} \emph{yādī}
  \item \emph{Say it:} \emph{yādī}, once more
  \item \emph{Say it:} \emph{jo tame yādī banāvśo to tame kaśũ nahī̃ bhūlo.} — If you make a list, you won't forget anything.
  \item \emph{Recall:} say the Gujarati for to throw down, then the Gujarati for to tire out (an opponent), then say \emph{yādī} again
\end{itemize}

\begin{tcolorbox}[breakable,colback=teal!4,colframe=teal!35!black,title={Before you move on}]
What does \gu{યાદી} mean? (\textquotedblleft{}A list\textquotedblright{}.) Say it once more.
\end{tcolorbox}
\section[\texorpdfstring{\gu{બજેટ} (bajeṭ)}{બજેટ (bajeṭ)}]{\gu{બજેટ} (bajeṭ) — a budget}

\label{lesson:GU-C528-bajet}



\begin{quote}
Before the new one: say the Gujarati for to tire out (an opponent), then the Gujarati for a list.
\end{quote}

\subsection*{You'll want to know: \gu{બજેટ}}
\textbf{\gu{બજેટ}} — \emph{bajeṭ} — \textquotedblleft{}a budget\textquotedblright{}.

\begin{quote}
\gu{જો} \gu{બજેટ} \gu{ઓછું} \gu{હશે} \gu{તો} \gu{આપણે} \gu{નાનો} \gu{કાર્યક્રમ} \gu{કરીશું}.
\end{quote}

\begin{quote}
\emph{jo bajeṭ ochũ haśe to āpaṇe nāno kāryakram karīśũ.}
\end{quote}

\begin{quote}
If the budget is small, we will hold a small event.
\end{quote}

A neuter noun.

\begin{cousinweb}
\textbf{\gu{બજેટ}} (\emph{bajeṭ}) is borrowed from English \emph{budget}.
\end{cousinweb}

\begin{grammarlens}[title={What you've built}]
One more word for saying what happens if.
\end{grammarlens}

\subsection*{Your turn}
\begin{itemize}
\raggedright
  \item \emph{Say it:} \emph{bajeṭ}
  \item \emph{Say it:} \emph{bajeṭ}, once more
  \item \emph{Say it:} \emph{jo bajeṭ ochũ haśe to āpaṇe nāno kāryakram karīśũ.} — If the budget is small, we will hold a small event.
  \item \emph{Recall:} say the Gujarati for to tire out (an opponent), then the Gujarati for a list, then say \emph{bajeṭ} again
\end{itemize}

\begin{tcolorbox}[breakable,colback=teal!4,colframe=teal!35!black,title={Before you move on}]
What does \gu{બજેટ} mean? (\textquotedblleft{}A budget\textquotedblright{}.) Say it once more.
\end{tcolorbox}
\section[\texorpdfstring{\gu{વ્યવસ્થા} (vyavasthā)}{વ્યવસ્થા (vyavasthā)}]{\gu{વ્યવસ્થા} (vyavasthā) — an arrangement}

\label{lesson:GU-C528-vyavastha}



\begin{quote}
Before the new one: say the Gujarati for a list, then the Gujarati for a budget.
\end{quote}

\subsection*{You'll want to know: \gu{વ્યવસ્થા}}
\textbf{\gu{વ્યવસ્થા}} — \emph{vyavasthā} — \textquotedblleft{}an arrangement\textquotedblright{}.

\begin{quote}
\gu{જો} \gu{મહેમાનો} \gu{વધારે} \gu{હશે} \gu{તો} \gu{આપણે} \gu{જમવાની} \gu{વ્યવસ્થા} \gu{બદલવી} \gu{પડશે}.
\end{quote}

\begin{quote}
\emph{jo mahemāno vadhāre haśe to āpaṇe jamvānī vyavasthā badalvī paḍśe.}
\end{quote}

\begin{quote}
If there are more guests, we will have to change the dinner arrangements.
\end{quote}

A feminine noun.

\begin{cousinweb}
\textbf{\gu{વ્યવસ્થા}} (\emph{vyavasthā}) is a tatsama, from Sanskrit \emph{vy-ava-sthā}, \textquotedblleft{}to stand apart, to be settled\textquotedblright{} (Monier-Williams): things set each in its place.
\end{cousinweb}

\begin{grammarlens}[title={What you've built}]
One more word for saying what happens if.
\end{grammarlens}

\subsection*{Your turn}
\begin{itemize}
\raggedright
  \item \emph{Say it:} \emph{vyavasthā}
  \item \emph{Say it:} \emph{vyavasthā}, once more
  \item \emph{Say it:} \emph{jo mahemāno vadhāre haśe to āpaṇe jamvānī vyavasthā badalvī paḍśe.} — If there are more guests, we will have to change the dinner arrangements.
  \item \emph{Recall:} say the Gujarati for a list, then the Gujarati for a budget, then say \emph{vyavasthā} again
\end{itemize}

\begin{tcolorbox}[breakable,colback=teal!4,colframe=teal!35!black,title={Before you move on}]
What does \gu{વ્યવસ્થા} mean? (\textquotedblleft{}An arrangement\textquotedblright{}.) Say it once more.
\end{tcolorbox}
\section[\texorpdfstring{\gu{ડાયરી} (ḍāyrī)}{ડાયરી (ḍāyrī)}]{\gu{ડાયરી} (ḍāyrī) — a diary}

\label{lesson:GU-C528-dayri}



\begin{quote}
Before the new one: say the Gujarati for a budget, then the Gujarati for an arrangement.
\end{quote}

\subsection*{You'll want to know: \gu{ડાયરી}}
\textbf{\gu{ડાયરી}} — \emph{ḍāyrī} — \textquotedblleft{}a diary\textquotedblright{}.

\begin{quote}
\gu{જો} \gu{તમે} \gu{ડાયરીમાં} \gu{લખશો} \gu{તો} \gu{તમને} \gu{બધી} \gu{તારીખો} \gu{યાદ} \gu{રહેશે}.
\end{quote}

\begin{quote}
\emph{jo tame ḍāyrīmā̃ lakhśo to tamne badhī tārīkho yād raheśe.}
\end{quote}

\begin{quote}
If you write it in a diary, you will remember all the dates.
\end{quote}

A feminine noun.

\begin{cousinweb}
\textbf{\gu{ડાયરી}} (\emph{ḍāyrī}) is borrowed from English \emph{diary}.
\end{cousinweb}

\begin{grammarlens}[title={What you've built}]
One more word for saying what happens if.
\end{grammarlens}

\subsection*{Your turn}
\begin{itemize}
\raggedright
  \item \emph{Say it:} \emph{ḍāyrī}
  \item \emph{Say it:} \emph{ḍāyrī}, once more
  \item \emph{Say it:} \emph{jo tame ḍāyrīmā̃ lakhśo to tamne badhī tārīkho yād raheśe.} — If you write it in a diary, you will remember all the dates.
  \item \emph{Recall:} say the Gujarati for a budget, then the Gujarati for an arrangement, then say \emph{ḍāyrī} again
\end{itemize}

\begin{tcolorbox}[breakable,colback=teal!4,colframe=teal!35!black,title={Before you move on}]
What does \gu{ડાયરી} mean? (\textquotedblleft{}A diary\textquotedblright{}.) Say it once more.
\end{tcolorbox}
\section[\texorpdfstring{\gu{યોજવું} (yojvũ)}{યોજવું (yojvũ)}]{\gu{યોજવું} (yojvũ) — to organise (an event)}

\label{lesson:GU-C528-yojvu}



\begin{quote}
Before the new one: say the Gujarati for an arrangement, then the Gujarati for a diary.
\end{quote}

\subsection*{You'll want to know: \gu{યોજવું}}
\textbf{\gu{યોજવું}} — \emph{yojvũ} — \textquotedblleft{}to organise (an event)\textquotedblright{}.

\begin{quote}
\gu{જો} \gu{વરસાદ} \gu{નહીં} \gu{પડે} \gu{તો} \gu{અમે} \gu{બગીચામાં} \gu{મેળો} \gu{યોજીશું}.
\end{quote}

\begin{quote}
\emph{jo varsād nahī̃ paḍe to ame bagīcāmā̃ meḷo yojīśũ.}
\end{quote}

\begin{quote}
If it doesn't rain, we will hold a fair in the garden.
\end{quote}

Past, with \gu{તેણે} (he or she …): \gu{તેણે} \gu{મેળો} \gu{યોજ્યો}, \emph{teṇe meḷo yojyo}.

\begin{cousinweb}
\textbf{\gu{યોજવું}} (\emph{yojvũ}) is a learned verb built on Sanskrit \emph{yojana}, \textquotedblleft{}joining, arranging\textquotedblright{} (Monier-Williams), the word behind \textbf{\gu{યોજના}}, a plan.
\end{cousinweb}

\begin{grammarlens}[title={What you've built}]
One more word for saying what happens if.
\end{grammarlens}

\subsection*{Your turn}
\begin{itemize}
\raggedright
  \item \emph{Say it:} \emph{yojvũ}
  \item \emph{Say it:} \emph{yojvũ}, once more
  \item \emph{Say it:} \emph{jo varsād nahī̃ paḍe to ame bagīcāmā̃ meḷo yojīśũ.} — If it doesn't rain, we will hold a fair in the garden.
  \item \emph{Recall:} say the Gujarati for an arrangement, then the Gujarati for a diary, then say \emph{yojvũ} again
\end{itemize}

\begin{tcolorbox}[breakable,colback=teal!4,colframe=teal!35!black,title={Before you move on}]
What does \gu{યોજવું} mean? (\textquotedblleft{}To organise (an event)\textquotedblright{}.) Say it once more.
\end{tcolorbox}
